\begin{solapartat}
\punts{0,5} Dibuix. El camp magnètic terrestre està dirigit en la direcció del fil i sentit nord de la brúixola, com es dedueix del circuit obert. El camp magnètic generat per la intensitat de corrent que passa pel fil conductor està dirigit cap a la dreta ja que l'agulla gira cap aquest costat.
\figura[0.3]{sol_fig1.png}

\punts{0,75} L'agulla és a sobre del fil conductor, ja que segons la direcció del corrent elèctric i utilitzant la regla de la mà dreta (o equivalent) amb aquesta distribució es genera un camp magnètic apuntant cap a la dreta en punts per sobre el corrent elèctric.
\end{solapartat}

\begin{solapartat}
\punts{0,5} Veiem que les components del camp magnètic total corresponen al camp magnètic terrestre i al camp generat pel fil conductor. Si l'agulla forma un angle de $45^\circ$ amb el corrent elèctric, els mòduls dels dos camps són iguals ja que $tg\,45^\circ = 1 = \frac{B_{fil}}{B_{Terra}}$ i per tant deduïm que $B_{Terra} = B_{fil}$.

\punts{0,5} Calculem el camp generat per la intensitat de corrent:
\[ B_{fil} = \frac{\mu_0 I}{2\pi r} = \frac{4\pi\cdot 10^{-7}\cdot 20}{2\pi\cdot 0,1} = 4,0\cdot 10^{-5}\ \mathrm{T} = 40\ \mathrm{\mu T} \]

\punts{0,25} Per tant, la component horitzontal del camp magnètic terrestre és $B_{Terra} = 40\ \mathrm{\mu T}$
\end{solapartat}
